\documentclass[11pt,a4paper]{article}

% 1. Paquetes de Idioma y Codificación
\usepackage[utf8]{inputenc}
\usepackage[T1]{fontenc}
\usepackage[spanish, es-noshorthands]{babel}

% 2. Paquetes de Diseño y Color (IMPORTANTE: xcolor con opción table)
\usepackage[table]{xcolor} 
\usepackage[table, x11names]{xcolor}
\usepackage{geometry}
\geometry{margin=2.5cm}
\usepackage{ragged2e}

% 3. Paquetes de Tablas y Listas
\usepackage{float}
\usepackage{array}
\usepackage{tabularx}
\usepackage{booktabs}
\usepackage{enumitem}

% 4. Definición de tipos de columna (Después de cargar tabularx y ragged2e)
\newcolumntype{L}{>{\RaggedRight\arraybackslash}X}

% 5. Otros paquetes
\usepackage{titlesec}
\usepackage{fancyhdr}
\usepackage{underscore} 

% --- Ajustes de cabecera ---
\setlength{\headheight}{14pt}
\addtolength{\topmargin}{-2pt}

% --- Estética ---
\titleformat{\section}{\large\bfseries\color{blue!70!black}}{}{0em}{\thesection.~}
\pagestyle{fancy}
\fancyhf{}
\fancyhead[L]{Sabor del Himalaya - PGCS}
\fancyhead[R]{Ciclo 2}
\fancyfoot[C]{\thepage}

\title{\textbf{Plan de Gestión de la Configuración Software}}
\author{Equipo 6}
\date{\today}

\begin{document}
\maketitle

% RECUERDA: Borra los archivos .toc y .aux antes de compilar
\tableofcontents 
\newpage

\section{Introducción}
El presente documento define la estrategia de Gestión de la Configuración Software (GCS) para el sistema de gestión del restaurante "Sabor del Himalaya". 

\section{Identificación de la Configuración}
\sloppy
Esta fase consiste en identificar y documentar los Elementos de Configuración (ECs) que forman parte del sistema.

\begin{itemize}
    \item \textbf{Código Backend (Propietario: Guillermo Guhl):}
        \begin{itemize}
            \item \texttt{SH-SRC-BE-APP}: Punto de entrada del servidor (\texttt{app.py}).
            \item \texttt{SH-SRC-BE-MODELS}: Definición de entidades de dominio (\texttt{models.py}).
            \item \texttt{SH-SRC-BE-REPOS}: Persistencia en MongoDB (\texttt{mongo\_repos.py}).
            \item \texttt{SH-SRC-BE-LOGIC}: Casos de uso de negocio (\texttt{order\_use\_cases.py}, \texttt{menu\_use\_cases.py}, \texttt{dish\_use\_cases.py}).
            \item \texttt{SH-SRC-BE-API}: Rutas y controladores REST (\texttt{order\_routes.py}, \texttt{menu\_routes.py}, \texttt{dish\_routes.py}).
        \end{itemize}
        
    \item \textbf{Código Frontend (Propietario: Daniel Gomarín):}
        \begin{itemize}
            \item \texttt{SH-SRC-FE-APP}: Componente principal de React (\texttt{App.jsx}).
            \item \texttt{SH-SRC-FE-VIEWS}: Vistas de navegación (\texttt{CartaView.jsx}, \texttt{DailyMenuView.jsx}).
            \item \texttt{SH-SRC-FE-COMP}: Componentes reutilizables (\texttt{DishCard.jsx}, \texttt{ShoppingCart.jsx}).
        \end{itemize}

    \item \textbf{Código Tests (Propietario: Carlos Sobrino):} 
        \begin{itemize}
            \item \texttt{SH-SRC-TEST}: Archivos de testing (\texttt{test\_billing\_buffer.py}, \texttt{test\_dish\_use\_cases.py}).
        \end{itemize}
    \item \textbf{Infraestructura y Entorno (Propietario: Marcos Romero):} 
        \begin{itemize}
            \item \texttt{SH-INF-DOCKER-COMP}: Orquestación de servicios (\texttt{docker-compose.yml}).
            \item \texttt{SH-INF-DOCKER-FILE}: Definición de imagen backend (\texttt{Dockerfile}).
        \end{itemize}

    \item \textbf{Documentación de Gestión (Propietario: Javier Pozuelo):}
        \begin{itemize}
            \item \texttt{SH-DOC-CV}: Ciclo de Vida del Proyecto.
            \item \texttt{SH-DOC-EE}: Estimación y Estrategia del Proyecto.
            \item \texttt{SH-DOC-C1}: Planificación del ciclo 1 del Proyecto.
            \item \texttt{SH-DOC-M1}: Monitorización del ciclo 1 del Proyecto.
            \item \texttt{SH-DOC-P1}: Presentación del ciclo 1 del Proyecto.
            \item \texttt{SH-DOC-C2}: Planificación del ciclo 2 del Proyecto.
            \item \texttt{SH-DOC-M2}: Monitorización del ciclo 2 del Proyecto.
            \item \texttt{SH-DOC-GC}: Gestión de la Calidad del Proyecto.
            \item \texttt{SH-DOC-P2}: Presentación del ciclo 2 del Proyecto.
            \item \texttt{SH-DOC-PGCS}: Gestión de la Configuración Software.
        \end{itemize}
\end{itemize}

El versionado seguirá el estándar \textbf{SemVer (Major.Minor.Patch)}.

\subsection{Definición de la línea base del sistema}

Para el proyecto "Sabor del Himalaya", se han establecido las siguientes líneas base:

\begin{itemize}
    \item \textbf{Línea Base 1 (Ciclo 1):} Corresponde a la versión oficial que satisface los requisitos \textbf{REQ 1} (Gestión de Platos), \textbf{REQ 2} (Menús Diarios), \textbf{REQ 3} (Navegación) y \textbf{REQ 9} (Login del Chef). Se consolidó tras el Ciclo 1.

    \item \textbf{Línea Base 2 (Ciclo 2):} Incorpora la implementación de la tramitación de pedidos (\textbf{REQ 4}) y el seguimiento de facturación (\textbf{REQ 8}). Se consolidará al final del ciclo 2.
\end{itemize}



\subsection{Procedimientos de copias de seguridad}

Para garantizar la integridad y la capacidad de recuperación del sistema ante fallos, se han establecido los siguientes procedimientos de salvaguarda, cumpliendo con el objetivo de mantener copias de todas las versiones para permitir la "marcha atrás" en caso de error:
\begin{itemize}
    \item \textbf{Repositorio Centralizado:} Se utiliza GitHub como repositorio principal para todos los Elementos de Configuración (ECs), incluyendo código fuente, documentos de diseño y planes de gestión.
    \item \textbf{Historial de Commits:} Cada \textit{commit} realizado en el repositorio actúa como una copia de seguridad incremental, registrando quién hizo el cambio, cuándo y qué se modificó, proporcionando un registro histórico completo.
    \item \textbf{Etiquetado de Líneas Base:} Siguiendo las directrices de GCS, cada vez que un conjunto de ECs se incorpora a una línea base oficial (Ciclo 1 o Ciclo 2), se creará un commit con una etiqueta fechada. Esto asegura que las versiones oficiales estén "etiquetadas y fechadas adecuadamente" para ser encontradas y recreadas fácilmente.
\end{itemize}

\subsection{Tabla Resumen de la Configuración}
\begin{table}[H]
\centering
\small
\begin{tabularx}{\textwidth}{@{}l l l l L@{}}
\toprule
\textbf{Nombre de EC} & \textbf{Acrónimo} & \textbf{Propietario} & \textbf{Fichero} & \textbf{Copia de Seguridad} \\ \midrule

\rowcolor{gray!15} \multicolumn{5}{l}{\textbf{Código Fuente - Backend}} \\
Punto entrada servidor & APP-BE & Guillermo Guhl & SH-SRC-BE-APP & GitHub \\
Entidades de dominio & MODELS & Guillermo Guhl & SH-SRC-BE-MODELS & GitHub \\
Persistencia MongoDB & REPOS & Guillermo Guhl & SH-SRC-BE-REPOS & GitHub \\
Lógica de negocio & LOGIC & Guillermo Guhl & SH-SRC-BE-LOGIC &  GitHub \\
Rutas y API REST & API & Guillermo Guhl & SH-SRC-BE-API & GitHub \\

\midrule
\rowcolor{gray!15} \multicolumn{5}{l}{\textbf{Código Fuente - Frontend}} \\
Componente React Principal & APP-FE & Daniel Gomarín & SH-SRC-FE-APP & GitHub \\
Vistas de navegación & VIEWS & Daniel Gomarín & SH-SRC-FE-VIEWS & GitHub \\
Componentes Reutilizables & COMP & Daniel Gomarín & SH-SRC-FE-COMP & GitHub \\

\midrule
\rowcolor{gray!15} \multicolumn{5}{l}{\textbf{Infraestructura y Entorno}} \\
Orquestación Docker & D-COMP & Marcos Romero & SH-INF-DOCKER-COMP & GitHub \\
Definición Imagen BE & D-FILE & Marcos Romero & SH-INF-DOCKER-FILE & GitHub \\

\midrule
\rowcolor{gray!15} \multicolumn{5}{l}{\textbf{Código Fuente - Tests}} \\
Archivos de testing & TESTS & Carlos Sobrino & SH-SRC-TEST & GitHub \\
\midrule
\rowcolor{gray!15} \multicolumn{5}{l}{\textbf{Documentos de Gestión}} \\
Ciclo de Vida & CV & Javier Pozuelo & SH-DOC-CV & Estrategia / GitHub \\
Estimación y Estrategia & EE & Javier Pozuelo & SH-DOC-EE & Estrategia / GitHub \\
Planificación Ciclo 1 & C1 & Javier Pozuelo & SH-DOC-C1 & Planif. / GitHub \\
Monitorización Ciclo 1 & M1 & Javier Pozuelo & SH-DOC-M1 & Seguim. / GitHub \\
Presentación Ciclo 1 & P1 & Javier Pozuelo & SH-DOC-P1 & Cierre / GitHub \\
Planificación Ciclo 2 & C2 & Javier Pozuelo & SH-DOC-C2 & Planif. / GitHub \\
Monitorización Ciclo 2 & M2 & Javier Pozuelo & SH-DOC-M2 & Seguim. / GitHub \\
Presentación Ciclo 2 & P2 & Javier Pozuelo & SH-DOC-P2 & Cierre / GitHub \\
Gestión de la Calidad & GC & Javier Pozuelo & SH-DOC-GC & Estrategia / GitHub \\
Plan de Gestión Conf. & PGCS & Javier Pozuelo & SH-DOC-PGCS & Estrategia / GitHub \\
\bottomrule
\end{tabularx}
\caption{Resumen completo de Elementos de Configuración y su ubicación.}
\end{table}

\section{Control de la Configuración}

\subsection{Comité de Control de la Configuración (CCC)}

El Comité de Control de la Configuración es el órgano responsable de supervisar la integridad del proyecto y validar cualquier cambio propuesto sobre la línea base. Para este proyecto, se ha definido una estructura que incluye los roles obligatorios de gestión y se ha optado por incorporar integrantes adicionales para reforzar las tareas de revisión y auditoría, garantizando así un control exhaustivo sobre el software.

\begin{table}[H]
\centering
\small
\begin{tabularx}{\textwidth}{@{}l l L@{}}
\toprule
\textbf{Integrante} & \textbf{Rol Principal} & \textbf{Funciones y Responsabilidades} \\ \midrule
Daniel Gomarín Ruesga & Responsable de Soporte & Gestiona el repositorio y las copias de seguridad. \\
Marcos Romero Villodres & Resp. de Desarrollo & Implementa los cambios propuestos. \\
Carlos Sobrino Martín & Resp. de Calidad & Monitoriza que el proceso de gestión de la configuración se realice según lo planificado. \\
Guillermo Guhl Arruabarrena & Integrante del CCC & Aprueba o rechaza los cambios propuestos mediante un proceso de revisión. \\
Javier Pozuelo Ruíz & Integrante del CCC & Aprueba o rechaza los cambios propuestos mediante un proceso de revisión. \\
\bottomrule
\end{tabularx}
\caption{Miembros y Roles del Comité de Control de la Configuración.}
\end{table}

\subsection{Peticiones de cambios a la configuración}

Durante el transcurso del proyecto, se han generado las siguientes peticiones de cambio oficiales para adaptar la planificación y la estructura del equipo a las necesidades del Ciclo 2:

\begin{itemize}
    \item \textbf{PC-002 (23/04/2026):} Asignación de roles Gestión de la Configuración y Gestión de la Calidad para cada integrante del equipo.
    \item \textbf{PC-003 (23/04/2026):} Modificación del alcance técnico eliminando el REQ 6 y sustituyéndolo por el REQ 8 para optimizar el tiempo disponible del equipo.
    \item \textbf{PC-004 (29/04/2026):} Modificación de la estructura de la base de datos, añadiendo dos nuevas colecciones para el cumplimiento de los REQs 4 y 8.
    \item \textbf{PC-005 (06/05/2026):} Incremento de la iluminación de las imágenes de los platos en la carta para su correcta visualización en la app.
\end{itemize}

\begin{table}[H]
\centering
\small
\begin{tabularx}{\textwidth}{@{}l l l L@{}}
\toprule
\textbf{ID PC} & \textbf{Fecha} & \textbf{Peticionario} & \textbf{Descripción del Cambio} \\ \midrule
PC-002 & 23/04/2026 & Marcos Romero & Asignación formal de roles GCS y GC. \\
PC-003 & 23/04/2026 & Guillermo Guhl & Sustitución de REQ 6 por REQ 8 por gestión de tiempos. \\
PC-004 & 29/04/2026 & Carlos Sobrino & Adición de nuevas colecciones en la base de datos. \\
PC-005 & 06/05/2026 & Javier Pozuelo & Aumento de iluminación de imágenes en la carta. \\
\bottomrule
\end{tabularx}
\caption{Resumen de Peticiones de Cambio generadas.}
\end{table}


\section{Informes de estado de configuración}

Esta sección detalla el seguimiento quincenal del estado de la configuración del sistema, permitiendo evaluar la estabilidad del producto y la eficacia de los procesos de control a lo largo del tiempo.

\subsection{Resumen Consolidado de Informes (Semanas 8-10)}
La siguiente tabla resume la evolución del proyecto basándose en los ccinco informes periódicos realizados.

\begin{table}[H]
\centering
\small
\begin{tabularx}{\textwidth}{@{}c c c c c L@{}}
\toprule
\textbf{Semana} & \textbf{PCs (Acum)} & \textbf{Págs (Acum)} & \textbf{LOC Totales} & \textbf{Actividad Relevante} \\ \midrule
8 & 3 & 72 & 1911 & Finalización primera versión (Ciclo 1) y PC-002, PC-003. \\
10 & 5 & 107 & 2279 & Finalización segunda versión (Ciclo 2) y PC-004, PC-005. \\
\bottomrule
\end{tabularx}
\caption{Evolución quincenal de la actividad y volumen de GCS.}
\end{table}

Se adjunta información más detallada en los informes.

\subsection{Historial de versiones}

Este historial registra la evolución oficial de la configuración del sistema, documentando únicamente los cambios en la línea base que han sido aprobados y cerrados por el Comité de Control de la Configuración.

\begin{table}[H]
\centering
\small
\begin{tabularx}{\textwidth}{@{}l l l L@{}}
\toprule
\textbf{Versión} & \textbf{Fecha} & \textbf{Referencia} & \textbf{Hito / Modificación} \\ \midrule
v1.0.0 & 27/03/2026 & Línea Base 1 & Consolidación oficial del Ciclo 1 tras la implementación de los requisitos REQ 1, 2, 3 y 9. \\
v1.1.0 & 29/04/2026 & PC-004 & Actualización del modelo de datos con nuevas colecciones para pedidos y facturación (REQs 4 y 8). \\
v1.2.0 & 06/05/2026 & PC-005 & Optimización de visualización en componentes Frontend para corregir problemas de iluminación. \\
v2.0.0 & 11/05/2026 & Línea Base 2 & Consolidación oficial del Ciclo 2 tras la realización de tests. \\
\bottomrule
\end{tabularx}
\caption{Historial de versiones oficiales del proyecto.}
\end{table}


\end{document}